\documentclass[12pt, a4paper]{article}
\usepackage[utf8]{inputenc}
\usepackage[T1]{fontenc}
\usepackage[french]{babel}
\usepackage[margin=2.5cm]{geometry}
\usepackage{amsmath, amssymb}
\usepackage{listings}
\usepackage{xcolor}
\usepackage{tcolorbox}
\tcbuselibrary{theorems, skins, breakable}
\usepackage{booktabs}
\usepackage{fancyhdr}
\usepackage{hyperref}
\usepackage{enumitem}

\definecolor{suva_blue}{RGB}{30, 100, 180}
\definecolor{code_bg}{RGB}{245, 247, 250}
\definecolor{code_kw}{RGB}{0, 80, 200}
\definecolor{code_str}{RGB}{180, 40, 40}
\definecolor{code_cmt}{RGB}{80, 130, 80}
\definecolor{warn_bg}{RGB}{255, 248, 220}
\definecolor{success_bg}{RGB}{230, 255, 235}
\definecolor{sol_green}{RGB}{0, 120, 80}
\definecolor{sol_bg}{RGB}{240, 255, 245}

\lstdefinestyle{python}{
  language=Python,
  basicstyle=\ttfamily\small,
  keywordstyle=\color{code_kw}\bfseries,
  stringstyle=\color{code_str},
  commentstyle=\color{code_cmt}\itshape,
  backgroundcolor=\color{code_bg},
  frame=single, rulecolor=\color{gray!40},
  numbers=left, numberstyle=\tiny\color{gray},
  breaklines=true, tabsize=4, showstringspaces=false,
}
\lstdefinestyle{solidity}{
  language=C++,
  basicstyle=\ttfamily\small,
  keywordstyle=\color{code_kw}\bfseries,
  stringstyle=\color{code_str},
  commentstyle=\color{code_cmt}\itshape,
  backgroundcolor=\color{sol_bg},
  frame=single, rulecolor=\color{sol_green!40},
  numbers=left, numberstyle=\tiny\color{gray},
  breaklines=true, tabsize=4, showstringspaces=false,
  morekeywords={uint256, uint8, bytes, bytes32, bool, address, memory,
    storage, external, view, pure, returns, mapping, struct,
    event, emit, require, assembly, let, mstore, mload},
}
\lstdefinestyle{bash}{
  basicstyle=\ttfamily\small\color{white},
  backgroundcolor=\color{gray!80},
  frame=single, rulecolor=\color{gray!80},
  breaklines=true,
}
\lstdefinestyle{output}{
  basicstyle=\ttfamily\small,
  backgroundcolor=\color{gray!10},
  frame=single, rulecolor=\color{gray!50},
  breaklines=true,
}

\newtcolorbox{exercice}[2][]{colback=white,colframe=suva_blue!60,
  fonttitle=\bfseries,title={Exercice #2},#1,breakable}
\newtcolorbox{solution}[1][]{colback=gray!5,colframe=gray!40,
  fonttitle=\bfseries\itshape,title={Solution},#1,breakable}
\newtcolorbox{objetif}[1][]{colback=success_bg,colframe=sol_green,
  fonttitle=\bfseries,title={Objectif},#1}
\newtcolorbox{warningbox}[1][]{colback=warn_bg,colframe=orange!80!black,
  fonttitle=\bfseries,title={Attention},#1}
\newtcolorbox{tipbox}[1][]{colback=suva_blue!5,colframe=suva_blue!50,
  fonttitle=\bfseries,title={Conseil},#1}

\pagestyle{fancy}
\fancyhf{}
\fancyhead[L]{\color{suva_blue}\textbf{SuVa --- Chapitre 8}}
\fancyhead[R]{\color{gray}Pratique Code}
\fancyfoot[L]{\footnotesize\color{gray}Mohamed Lamine MBENGUE}
\fancyfoot[C]{\thepage}
\fancyfoot[R]{\footnotesize\color{gray}portfolio-alamine.web.app}
\renewcommand{\headrulewidth}{0.4pt}
\renewcommand{\footrulewidth}{0.3pt}

% ══════════════════════════════════════════════════════════════════════
\begin{document}

\begin{center}
  {\color{suva_blue}\rule{\linewidth}{2pt}}\\[0.5cm]
  {\Huge\bfseries\color{suva_blue} Chapitre 8}\\[0.3cm]
  {\large\bfseries Pratique Code --- Maîtriser le projet par l'action}\\[0.2cm]
  {\small\color{gray} Exercices directement sur les fichiers du projet ETHDILITHIUM}\\[0.3cm]
  {\color{suva_blue}\rule{\linewidth}{2pt}}
\end{center}

\vspace{0.3cm}
\begin{tcolorbox}[colback=suva_blue!5, colframe=suva_blue!40,
  left=6pt, right=6pt, top=4pt, bottom=4pt]
\small
\begin{tabular}{ll}
  \textbf{Auteur}      & Mohamed Lamine MBENGUE \\
  \textbf{Spécialités} & Sécurité · DevOps · Dev Full Stack Web \& Mobile \\
  \textbf{Contact}     & +221 78 178 44 36 · mbenguemohamedlamine2022@gmail.com \\
  \textbf{Portfolio}   & \href{https://portfolio-alamine.web.app}{\texttt{portfolio-alamine.web.app}} \\
\end{tabular}
\end{tcolorbox}
\vspace{0.3cm}

\tableofcontents
\newpage

%══════════════════════════════════════════════════════════════════════
\section{Structure du projet --- Vue d'ensemble}
%══════════════════════════════════════════════════════════════════════

Avant de coder, connais la carte du terrain :

\begin{tcolorbox}[colback=gray!5, colframe=gray!50, breakable]
\begin{verbatim}
ETHDILITHIUM-main/
├── makefile                    # make install / make test / make gen_test_vectors
├── python-ref/
│   ├── makefile                # make install (pip) / make test (pytest)
│   ├── requirements.txt        # pycryptodome, eth_abi
│   └── Falcon_dilithium_py/
│       ├── Falcon_dilithium/
│       │   ├── Falcon_dilithium.py         # CLASSE PRINCIPALE : keygen, sign, verify
│       │   ├── Falcon_dilithium_eth.py     # Variante ETHDilithium (Keccak PRNG)
│       │   └── Falcon_default_parameters.py # Paramètres Dilithium2/3/5
│       ├── utilities/
│       │   └── utils.py                    # decompose, HighBits, LowBits, MakeHint, UseHint
│       ├── polynomials/
│       │   └── polynomials.py              # PolynomialRingDilithium (NTT-based mult)
│       ├── modules/
│       │   └── modules.py                  # ModuleDilithium (matrix-vector A@s)
│       ├── Falcon_keccak_prng/
│       │   └── keccak_prng_wrapper.py      # PRNG Keccak compatible Solidity
│       ├── shake/
│       │   └── shake_wrapper.py            # SHAKE-128/256
│       └── tests/
│           └── test_dilithium.py           # pytest : KAT vectors
└── assets/
    └── PQCsignKAT_Dilithium2.rsp           # Vecteurs de test NIST
\end{verbatim}
\end{tcolorbox}

\begin{objetif}
Après ce chapitre tu dois pouvoir :
\begin{itemize}
  \item Naviguer dans le projet sans hésiter
  \item Lire et modifier \texttt{Falcon\_dilithium.py} en comprenant chaque ligne
  \item Écrire tes propres tests Python et les faire passer
  \item Comprendre les paramètres et leur impact sur la sécurité
  \item Tracer le flux complet keygen $\to$ sign $\to$ verify
\end{itemize}
\end{objetif}

%══════════════════════════════════════════════════════════════════════
\section{Mise en place de l'environnement}
%══════════════════════════════════════════════════════════════════════

\subsection{Installation complète (une seule fois)}

\begin{lstlisting}[style=bash]
# Depuis le dossier ETHDILITHIUM-main/
cd python-ref

# Créer l'environnement virtuel
python -m venv myenv

# Activer (Linux/Mac)
source myenv/bin/activate
# Activer (Windows PowerShell)
.\myenv\Scripts\Activate.ps1
# Activer (Windows Git Bash)
source myenv/Scripts/activate

# Installer les dépendances
pip install -r requirements.txt

# Installer le package polyntt (NTT optimisé de ZKNox)
pip install git+https://github.com/ZKNox/polyntt.git

# Installer le package dilithium_py en mode développement
pip install -e .
\end{lstlisting}

\subsection{Lancer les tests existants}

\begin{lstlisting}[style=bash]
# Depuis python-ref/Falcon_dilithium_py/
cd Falcon_dilithium_py

# Lancer tous les tests
python -m pytest tests/ -v

# Lancer un test spécifique
python -m pytest tests/test_dilithium.py::test_dilithium2 -v

# Avec affichage des prints
python -m pytest tests/ -v -s
\end{lstlisting}

\begin{lstlisting}[style=output]
tests/test_dilithium.py::test_dilithium2 PASSED        [ 25%]
tests/test_dilithium.py::test_dilithium3 PASSED        [ 50%]
tests/test_dilithium.py::test_dilithium5 PASSED        [ 75%]
tests/test_dilithium.py::test_ethdilithium2_keccak_prng PASSED [100%]

4 passed in 12.34s
\end{lstlisting}

\newpage
%══════════════════════════════════════════════════════════════════════
\section{Exercices --- Partie 1 : Exploration des classes}
%══════════════════════════════════════════════════════════════════════

%──────────────────────────────────────────────────────────────────────
\subsection{Exercice P1 : Premier keygen/sign/verify}
%──────────────────────────────────────────────────────────────────────

\begin{exercice}{P1 --- Ton premier keygen-sign-verify}
Crée un fichier \texttt{mon\_test.py} dans le dossier \texttt{Falcon\_dilithium\_py/} et écris un script qui :
\begin{enumerate}
  \item Importe la classe \texttt{Dilithium2}
  \item Génère une paire de clés
  \item Signe le message \texttt{"SuVa Protocol - quantum proof!"}
  \item Vérifie la signature
  \item Affiche la taille de chaque objet (pk, sk, sig)
  \item Essaie de vérifier avec un message modifié (doit échouer)
\end{enumerate}
\end{exercice}

\begin{solution}
\begin{lstlisting}[style=python]
# Falcon_dilithium_py/mon_test.py
import sys
sys.path.insert(0, '.')  # S'assurer que le module est trouvé

from Falcon_dilithium.Falcon_dilithium import Dilithium2

# 1. Instancier Dilithium2
d = Dilithium2

# 2. Génération des clés
pk, sk = d.keygen()
print(f"[+] KeyGen OK")
print(f"    Public key  : {len(pk)} bytes")
print(f"    Secret key  : {len(sk)} bytes")

# 3. Signature
message = b"SuVa Protocol - quantum proof!"
sig = d.sign(sk, message)
print(f"\n[+] Signature OK")
print(f"    Signature   : {len(sig)} bytes")

# 4. Vérification
valid = d.verify(pk, message, sig)
print(f"\n[+] Vérification : {'OK (True)' if valid else 'ECHEC (False)'}")

# 5. Mauvais message
bad_message = b"SuVa Protocol - quantum proof?"  # '?' au lieu de '!'
invalid = d.verify(pk, bad_message, sig)
print(f"[+] Mauvais message : {'ECHEC (attendu)' if not invalid else 'PROBLEME!'}")

# Affichage hex des 16 premiers octets
print(f"\n--- Extrait public key ---")
print(f"    {pk[:16].hex()}...")
print(f"--- Extrait signature ---")
print(f"    {sig[:16].hex()}...")
\end{lstlisting}

\textbf{Pour exécuter :}
\begin{lstlisting}[style=bash]
python mon_test.py
\end{lstlisting}

\textbf{Sortie attendue :}
\begin{lstlisting}[style=output]
[+] KeyGen OK
    Public key  : 1312 bytes
    Secret key  : 2560 bytes

[+] Signature OK
    Signature   : 2420 bytes

[+] Vérification : OK (True)
[+] Mauvais message : ECHEC (attendu)
\end{lstlisting}
\end{solution}

%──────────────────────────────────────────────────────────────────────
\subsection{Exercice P2 : Explorer les paramètres}
%──────────────────────────────────────────────────────────────────────

\begin{exercice}{P2 --- Comparer les niveaux de sécurité}
Ouvre le fichier \texttt{Falcon\_default\_parameters.py} et écris un script qui :
\begin{enumerate}
  \item Affiche les paramètres clés de Dilithium2, Dilithium3, Dilithium5 dans un tableau
  \item Pour chaque niveau, génère une clé et une signature et affiche les tailles
  \item Mesure le temps de keygen et sign avec \texttt{time.time()}
\end{enumerate}

\textbf{Paramètres à afficher :} \texttt{n, q, k, l, eta, beta, gamma1, gamma2}
\end{exercice}

\begin{solution}
\begin{lstlisting}[style=python]
# Falcon_dilithium_py/compare_params.py
import time
from Falcon_dilithium.Falcon_dilithium import Dilithium2, Dilithium3, Dilithium5

schemes = [
    ("Dilithium2", Dilithium2),
    ("Dilithium3", Dilithium3),
    ("Dilithium5", Dilithium5),
]

print(f"{'Scheme':<15} {'n':>5} {'q':>10} {'k':>4} {'l':>4} "
      f"{'eta':>5} {'beta':>6} {'PK':>6} {'SK':>6} {'Sig':>6} "
      f"{'Keygen(ms)':>12} {'Sign(ms)':>10}")
print("-" * 95)

msg = b"test message"

for name, D in schemes:
    params = D.params
    n   = params['n']
    q   = params['q']
    k   = params['k']
    l   = params['l']
    eta = params['eta']
    beta = params['beta']

    # Timing keygen
    t0 = time.time()
    pk, sk = D.keygen()
    t_keygen = (time.time() - t0) * 1000

    # Timing sign
    t0 = time.time()
    sig = D.sign(sk, msg)
    t_sign = (time.time() - t0) * 1000

    print(f"{name:<15} {n:>5} {q:>10} {k:>4} {l:>4} "
          f"{eta:>5} {beta:>6} {len(pk):>6} {len(sk):>6} {len(sig):>6} "
          f"{t_keygen:>12.1f} {t_sign:>10.1f}")
\end{lstlisting}

\textbf{Sortie attendue :}
\begin{lstlisting}[style=output]
Scheme            n          q    k    l   eta   beta     PK     SK    Sig   Keygen(ms)   Sign(ms)
-----------------------------------------------------------------------------------------------
Dilithium2     256    8380417    4    4     2    78   1312   2560   2420       45.2       38.7
Dilithium3     256    8380417    6    5     4   196   1952   4000   3293       62.1       55.3
Dilithium5     256    8380417    8    7     2   120   2592   4864   4595       88.4       75.2
\end{lstlisting}
\end{solution}

\newpage
%══════════════════════════════════════════════════════════════════════
\section{Exercices --- Partie 2 : Modifier et comprendre l'algorithme}
%══════════════════════════════════════════════════════════════════════

%──────────────────────────────────────────────────────────────────────
\subsection{Exercice P3 : Observer le Rejection Sampling}
%──────────────────────────────────────────────────────────────────────

\begin{exercice}{P3 --- Compter les itérations du rejection sampling}
Le signing de Dilithium boucle tant que certaines conditions ne sont pas satisfaites (rejection sampling).
Modifie \texttt{Falcon\_dilithium.py} pour compter le nombre d'itérations.

\textbf{Lieu d'intervention :} méthode \texttt{sign()} --- la boucle \texttt{while True}

Signe 100 messages différents et affiche : min, max, moyenne des itérations.
\end{exercice}

\begin{solution}
Ouvre \texttt{Falcon\_dilithium/Falcon\_dilithium.py} et trouve la boucle \texttt{while True} dans \texttt{sign()}.
Ajoute un compteur :

\begin{lstlisting}[style=python]
# Dans la méthode sign(), trouver le while True:
# Ajouter AVANT la boucle :
iteration_count = 0

# Au DEBUT du corps de la boucle while True: ajouter :
iteration_count += 1

# Après le return sig (ou à la sortie de la boucle), ajouter :
# (NB: en Python le while True se termine par break ou return)
# Retourner aussi le compteur pour les stats (ou utiliser une variable globale)
\end{lstlisting}

\textbf{Script de test séparé :}
\begin{lstlisting}[style=python]
# Falcon_dilithium_py/count_iterations.py
import random
from Falcon_dilithium.Falcon_dilithium import Dilithium2

# Patcher temporairement pour compter
d = Dilithium2
pk, sk = d.keygen()

counts = []
N = 100
for i in range(N):
    msg = f"message_{i}_{random.randint(0, 2**32)}".encode()
    # Modifier temporairement sign() pour retourner (sig, count)
    # Alternative : mesurer le temps (corrélé aux itérations)
    import time
    t0 = time.time()
    sig = d.sign(sk, msg)
    elapsed = (time.time() - t0) * 1000
    counts.append(elapsed)
    assert d.verify(pk, msg, sig)

print(f"Signing {N} messages:")
print(f"  Min  : {min(counts):.1f} ms")
print(f"  Max  : {max(counts):.1f} ms")
print(f"  Mean : {sum(counts)/len(counts):.1f} ms")
print(f"\nNote: Le temps est proportionnel aux itérations.")
print(f"La moyenne théorique est ~4.25 itérations.")
\end{lstlisting}

\textbf{Modification directe dans Falcon\_dilithium.py :}
\begin{lstlisting}[style=python]
def sign(self, sk, m):
    # ... (code existant pour unpack sk) ...

    kappa = 0
    iteration = 0  # <-- AJOUTER ICI

    while True:
        iteration += 1  # <-- AJOUTER ICI
        y = self._expandMask(rhoprime, kappa)
        # ... reste du code ...
        if cs2_norm >= self.gamma2 - self.beta:
            kappa += self.l
            continue
        if hints_n > self.omega:
            kappa += self.l
            continue
        # La signature est valide
        # print(f"Iterations: {iteration}")  # <-- DÉCOMMENTER pour debug
        return self._pack_sig(...)
\end{lstlisting}
\end{solution}

%──────────────────────────────────────────────────────────────────────
\subsection{Exercice P4 : Étudier les fonctions utilitaires}
%──────────────────────────────────────────────────────────────────────

\begin{exercice}{P4 --- Tester decompose(), HighBits(), LowBits()}
Écris un script de test pour \texttt{utilities/utils.py}. Tu dois :
\begin{enumerate}
  \item Tester \texttt{decompose(r, alpha, q)} sur plusieurs valeurs
  \item Vérifier que \texttt{r == r1 * alpha + r0} (relation fondamentale)
  \item Tester \texttt{high\_bits()} et \texttt{low\_bits()} séparément
  \item Tester \texttt{check\_norm\_bound()} sur un vecteur de polynômes
\end{enumerate}
\end{exercice}

\begin{solution}
\begin{lstlisting}[style=python]
# Falcon_dilithium_py/test_utils.py
import sys
sys.path.insert(0, '.')

from utilities.utils import (
    decompose, high_bits, low_bits,
    make_hint, use_hint, check_norm_bound
)

q = 8380417

# Paramètres Dilithium2
gamma2 = (q - 1) // 32  # = 261888
alpha  = 2 * gamma2     # = 523776

print("=== Test decompose() ===")
test_values = [0, 1, gamma2, gamma2 + 1, q // 2, q - 1]

for r in test_values:
    r1, r0 = decompose(r, alpha, q)
    # Vérification : r = r1 * alpha + r0 (mod q), r0 dans [-alpha/2, alpha/2]
    reconstructed = (r1 * alpha + r0) % q
    ok = (reconstructed == r % q)
    print(f"  r={r:>8}  =>  r1={r1:>4}, r0={r0:>8}  | Vérifié: {ok}")

print("\n=== Test high_bits() / low_bits() ===")
for r in [100, gamma2 - 1, gamma2, gamma2 + 1, 2 * gamma2]:
    h = high_bits(r, alpha, q)
    l = low_bits(r, alpha, q)
    print(f"  r={r:>8}  =>  high={h:>4}, low={l:>8}")

print("\n=== Test check_norm_bound() ===")
# Importer une structure de polynôme
from polynomials.polynomials import PolynomialRingDilithium

R = PolynomialRingDilithium()
# Créer un polynôme avec des coefficients petits (gamma1 - beta = large)
beta = 78  # Dilithium2
gamma1 = 2**17  # 131072

# Polynôme dont les coeff sont tous = beta - 1 (OK)
p_ok = R([beta - 1] * 256)
# Polynôme dont un coeff dépasse la borne
p_bad = R([gamma1] * 256)

print(f"  p_ok (coeff={beta-1}) bound={gamma1 - beta}: "
      f"{'PASSE' if not check_norm_bound(p_ok, gamma1 - beta) else 'ECHOUE'}")
print(f"  p_bad (coeff={gamma1}) bound={gamma1 - beta}: "
      f"{'PASSE' if not check_norm_bound(p_bad, gamma1 - beta) else 'ECHOUE'}")
\end{lstlisting}
\end{solution}

\newpage
%══════════════════════════════════════════════════════════════════════
\section{Exercices --- Partie 3 : ETHDilithium et la variante Ethereum}
%══════════════════════════════════════════════════════════════════════

%──────────────────────────────────────────────────────────────────────
\subsection{Exercice P5 : Comparer Dilithium vs ETHDilithium}
%──────────────────────────────────────────────────────────────────────

\begin{exercice}{P5 --- Dilithium standard vs variante Ethereum}
Ouvre \texttt{Falcon\_dilithium\_eth.py} et compare avec \texttt{Falcon\_dilithium.py}.
\begin{enumerate}
  \item Identifie les 3 différences principales entre les deux implémentations
  \item Écris un script qui génère une clé ETHDilithium et signe un message
  \item Montre que les signatures ETHDilithium et Dilithium standard sont \textbf{incompatibles}
\end{enumerate}
\end{exercice}

\begin{solution}
\textbf{Les 3 différences principales :}

\begin{tcolorbox}[colback=suva_blue!5, colframe=suva_blue!40]
\begin{enumerate}
  \item \textbf{PRNG :} ETHDilithium utilise \textbf{Keccak-256} (natif EVM) au lieu de SHAKE-256
  \item \textbf{Clé publique :} ETHDilithium inclut $\hat{A}$ (matrix NTT précalculée) dans la pk pour économiser du gas on-chain
  \item \textbf{Hash :} ETHDilithium hash le message avec Keccak-256 (opcode EVM 3 gas) au lieu de SHAKE-256 (coûteux car implémenté en Solidity)
\end{enumerate}
\end{tcolorbox}

\begin{lstlisting}[style=python]
# Falcon_dilithium_py/test_eth_vs_standard.py
from Falcon_dilithium.Falcon_dilithium import Dilithium2
from Falcon_dilithium.Falcon_dilithium_eth import ETHDilithium2KeccakPRNG

msg = b"SuVa stablecoin"

# --- Dilithium2 standard ---
pk_std, sk_std = Dilithium2.keygen()
sig_std = Dilithium2.sign(sk_std, msg)
print(f"Dilithium2 standard:")
print(f"  PK: {len(pk_std)} bytes")
print(f"  SK: {len(sk_std)} bytes")
print(f"  Sig: {len(sig_std)} bytes")
print(f"  Verify: {Dilithium2.verify(pk_std, msg, sig_std)}")

print()

# --- ETHDilithium2 Keccak PRNG ---
pk_eth, sk_eth = ETHDilithium2KeccakPRNG.keygen()
sig_eth = ETHDilithium2KeccakPRNG.sign(sk_eth, msg)
print(f"ETHDilithium2 (Keccak PRNG):")
print(f"  PK: {len(pk_eth)} bytes  (plus grande : contient A_hat)")
print(f"  SK: {len(sk_eth)} bytes")
print(f"  Sig: {len(sig_eth)} bytes")
print(f"  Verify: {ETHDilithium2KeccakPRNG.verify(pk_eth, msg, sig_eth)}")

print()

# --- Incompatibilité ---
print("=== Test d'incompatibilité ===")
try:
    result = Dilithium2.verify(pk_eth, msg, sig_eth)
    print(f"  Dilithium2.verify(pk_eth, sig_eth) = {result}")
    # Devrait retourner False ou lever une exception
except Exception as e:
    print(f"  Exception (attendue): {type(e).__name__}: {e}")
\end{lstlisting}
\end{solution}

%──────────────────────────────────────────────────────────────────────
\subsection{Exercice P6 : Vecteurs de test NIST (KAT)}
%──────────────────────────────────────────────────────────────────────

\begin{exercice}{P6 --- Lire et utiliser les vecteurs KAT du NIST}
Le fichier \texttt{assets/PQCsignKAT\_Dilithium2.rsp} contient les vecteurs de test officiels du NIST.
\begin{enumerate}
  \item Ouvre ce fichier et comprends son format
  \item Écris un parser qui lit le premier vecteur (count=0)
  \item Vérifie que le code Python produit la même signature que le KAT
\end{enumerate}
\end{exercice}

\begin{solution}
\begin{lstlisting}[style=python]
# Falcon_dilithium_py/test_kat.py
"""
Format du fichier KAT :
# DILITHIUM2
count = 0
seed = <64 hex chars>
mlen = <int>
msg = <hex>
pk = <hex>
sk = <hex>
smlen = <int>
sm = <signature+message hex>
"""
import os

def parse_kat_file(filepath, count=0):
    """Parse un vecteur KAT par son numéro (count)."""
    vectors = {}
    current = {}
    target_found = False

    with open(filepath, 'r') as f:
        for line in f:
            line = line.strip()
            if not line or line.startswith('#'):
                continue
            if '=' in line:
                key, val = line.split('=', 1)
                key = key.strip()
                val = val.strip()

                if key == 'count':
                    if current:
                        vectors[current['count']] = current
                    current = {'count': int(val)}
                elif key in ['mlen', 'smlen']:
                    current[key] = int(val)
                else:
                    current[key] = bytes.fromhex(val)

    if current:
        vectors[current['count']] = current

    return vectors.get(count)

# Chemin vers les KAT
kat_path = os.path.join(
    '..', '..', 'assets', 'PQCsignKAT_Dilithium2.rsp'
)

kat = parse_kat_file(kat_path, count=0)
print(f"KAT count=0:")
print(f"  seed   : {kat['seed'].hex()[:32]}...")
print(f"  mlen   : {kat['mlen']}")
print(f"  pk     : {kat['pk'].hex()[:32]}...")
print(f"  sk     : {kat['sk'].hex()[:32]}...")
print(f"  smlen  : {kat['smlen']}")
print(f"  sig    : {kat['sm'][:16].hex()}...")

# La signature dans le fichier KAT est le dernier champ 'sm'
# qui contient sig || msg (selon le format NIST)
sig_from_kat = kat['sm'][:kat['smlen'] - kat['mlen']]
msg_from_kat = kat['sm'][kat['smlen'] - kat['mlen']:]

print(f"\nMessage extrait : {msg_from_kat.hex()}")
print(f"Signature taille : {len(sig_from_kat)} bytes")

# Vérifier avec notre implémentation
from Falcon_dilithium.Falcon_dilithium import Dilithium2
result = Dilithium2.verify(kat['pk'], msg_from_kat, sig_from_kat)
print(f"\nVérification KAT : {'PASSE' if result else 'ECHEC'}")
\end{lstlisting}
\end{solution}

\newpage
%══════════════════════════════════════════════════════════════════════
\section{Exercices --- Partie 4 : NTT et Polynômes}
%══════════════════════════════════════════════════════════════════════

%──────────────────────────────────────────────────────────────────────
\subsection{Exercice P7 : Multiplication de polynômes NTT vs naïve}
%──────────────────────────────────────────────────────────────────────

\begin{exercice}{P7 --- Comparer NTT et multiplication naïve}
Implémente une multiplication naïve de polynômes (O(n²)) et compare avec la NTT (O(n log n)) :
\begin{enumerate}
  \item Implémente \texttt{multiply\_naive(a, b, n, q)} en Python pur
  \item Utilise \texttt{PolynomialRingDilithium} pour la multiplication NTT
  \item Vérifie qu'ils donnent le même résultat
  \item Compare les temps pour n=256
\end{enumerate}
\end{exercice}

\begin{solution}
\begin{lstlisting}[style=python]
# Falcon_dilithium_py/test_ntt_vs_naive.py
import time
import random
from polynomials.polynomials import PolynomialRingDilithium

q = 8380417
n = 256

def multiply_naive(a_coeffs, b_coeffs, n, q):
    """Multiplication naïve dans R_q = Z_q[X]/(X^n + 1)"""
    result = [0] * n
    for i in range(n):
        for j in range(n):
            idx = (i + j) % n
            # Si i + j >= n, on est dans X^n = -1, donc signe -1
            sign = -1 if (i + j) >= n else 1
            result[idx] = (result[idx] + sign * a_coeffs[i] * b_coeffs[j]) % q
    return result

R = PolynomialRingDilithium()

# Générer des polynômes aléatoires
random.seed(42)
a_coeffs = [random.randint(0, q - 1) for _ in range(n)]
b_coeffs = [random.randint(0, q - 1) for _ in range(n)]

# Création des objets polynomial
a_poly = R(a_coeffs)
b_poly = R(b_coeffs)

# --- Multiplication naïve ---
t0 = time.time()
c_naive = multiply_naive(a_coeffs, b_coeffs, n, q)
t_naive = (time.time() - t0) * 1000

# --- Multiplication NTT ---
t0 = time.time()
c_ntt_poly = a_poly * b_poly
t_ntt = (time.time() - t0) * 1000
c_ntt = list(c_ntt_poly.coeffs)

# Vérification
if c_naive == c_ntt:
    print("[PASSE] Les deux méthodes donnent le même résultat!")
else:
    diff = sum(1 for x, y in zip(c_naive, c_ntt) if x != y)
    print(f"[ECHEC] {diff} coefficients différents!")

print(f"\nTemps multiplication naïve  : {t_naive:.2f} ms")
print(f"Temps multiplication NTT    : {t_ntt:.4f} ms")
print(f"Accélération                : {t_naive/t_ntt:.0f}x")
print(f"\nn={n} => naïve: O(n²)={n**2} ops, NTT: O(n·log2(n))={n*8} ops")
\end{lstlisting}
\end{solution}

%──────────────────────────────────────────────────────────────────────
\subsection{Exercice P8 : Implémenter NTT from scratch}
%──────────────────────────────────────────────────────────────────────

\begin{exercice}{P8 --- Implémenter la NTT négacyclique}
Implémente la NTT de Dilithium (négacyclique, q=8380417, ζ=1753) depuis zéro.

\textbf{Algorithme Cooley-Tukey pour la NTT négacyclique :}
\begin{enumerate}
  \item Calculer la liste des puissances de ζ : \texttt{zetas[i]} = $\zeta^{\text{bitrev}(i)} \mod q$ pour $i = 0..127$
  \item Implémenter le butterfly : \texttt{t = zeta * b; a, b = a + t, a - t (mod q)}
  \item Tester que \texttt{INTT(NTT(a)) == a}
\end{enumerate}
\end{exercice}

\begin{solution}
\begin{lstlisting}[style=python]
# Falcon_dilithium_py/my_ntt.py
"""
NTT négacyclique pour Dilithium
q = 8380417, n = 256, zeta = 1753 (générateur primitif 512e racine de l'unité)
"""

q     = 8380417
n     = 256
zeta  = 1753  # tel que zeta^256 = -1 mod q (512e racine primitive de l'unité)

def mod_pow(base, exp, mod):
    """Exponentiation rapide modulaire."""
    result = 1
    base %= mod
    while exp > 0:
        if exp % 2 == 1:
            result = result * base % mod
        base = base * base % mod
        exp //= 2
    return result

def bit_reverse(k, log2n):
    """Inverse les bits de k sur log2n bits."""
    result = 0
    for _ in range(log2n):
        result = (result << 1) | (k & 1)
        k >>= 1
    return result

# Précalcul des facteurs de torsion (zetas)
# zetas[i] = zeta^bitrev(i) mod q pour i = 0..127
log2n = 8  # log2(256) = 8, mais on a 128 twiddle factors (n/2)
zetas = [mod_pow(zeta, bit_reverse(i, log2n - 1), q) for i in range(n // 2)]

def ntt(a):
    """NTT négacyclique de Dilithium (in-place, Cooley-Tukey)."""
    a = list(a)
    k = 0
    length = n // 2  # 128
    while length >= 1:
        start = 0
        while start < n:
            k += 1
            zeta_k = zetas[k] if k < len(zetas) else 1
            for j in range(start, start + length):
                t = zeta_k * a[j + length] % q
                a[j + length] = (a[j] - t) % q
                a[j]          = (a[j] + t) % q
            start += 2 * length
        length //= 2
    return a

def intt(a):
    """INTT négacyclique (inverse de ntt)."""
    a = list(a)
    k = n // 2  # 128
    length = 1
    while length <= n // 2:
        start = 0
        while start < n:
            k -= 1
            zeta_k_inv = mod_pow(zetas[k] if k < len(zetas) else 1, q - 2, q)
            for j in range(start, start + length):
                t         = a[j]
                a[j]      = (t + a[j + length]) % q
                a[j + length] = (zeta_k_inv * (a[j + length] - t)) % q
            start += 2 * length
        length *= 2
    # Normaliser par n
    n_inv = mod_pow(n, q - 2, q)
    return [(x * n_inv) % q for x in a]

# Test : INTT(NTT(a)) == a
import random
random.seed(0)
a = [random.randint(0, q - 1) for _ in range(n)]
a_ntt  = ntt(a)
a_back = intt(a_ntt)

if a == a_back:
    print("[PASSE] INTT(NTT(a)) == a  ✓")
else:
    diffs = sum(1 for x, y in zip(a, a_back) if x != y)
    print(f"[ECHEC] {diffs} coefficients incorrects")

# Comparer avec la bibliothèque polyntt
from polynomials.polynomials import PolynomialRingDilithium
R = PolynomialRingDilithium()
p = R(a)
p_ntt_lib = list(p.ntt().coeffs)

# Vérifier que notre NTT correspond
same = (a_ntt == p_ntt_lib)
print(f"[{'PASSE' if same else 'ECHEC'}] Notre NTT == polyntt library: {same}")
\end{lstlisting}
\end{solution}

\newpage
%══════════════════════════════════════════════════════════════════════
\section{Exercices --- Partie 5 : Comprendre le code Solidity}
%══════════════════════════════════════════════════════════════════════

\begin{warningbox}
Pour ces exercices, tu as besoin de \textbf{Foundry} installé.
\begin{lstlisting}[style=bash]
# Installation Foundry
curl -L https://foundry.paradigm.xyz | bash
foundryup

# Compiler le projet
cd ETHDILITHIUM-main
forge build

# Lancer les tests Solidity
forge test -vvv --gas-report
\end{lstlisting}
\end{warningbox}

%──────────────────────────────────────────────────────────────────────
\subsection{Exercice P9 : Lire ZKNOX\_NTT.sol et comprendre le Yul}
%──────────────────────────────────────────────────────────────────────

\begin{exercice}{P9 --- Analyser ZKNOX\_NTT.sol}
Ouvre \texttt{ETHDILITHIUM-main/src/ZKNOX\_NTT.sol} et réponds :
\begin{enumerate}
  \item Pourquoi utilise-t-on \texttt{extcodecopy} pour lire les twiddle factors ?
  \item Que fait l'instruction \texttt{assembly \{ ... \}} dans le butterfly ?
  \item Combien de twiddle factors sont stockés et comment ?
  \item Quel est le coût gas estimé de la NTT complète ?
\end{enumerate}
\end{exercice}

\begin{solution}
\begin{tcolorbox}[colback=suva_blue!5, colframe=suva_blue!40, breakable]
\textbf{1. Pourquoi \texttt{extcodecopy} ?}\\
Les 128 twiddle factors ($\zeta^{\text{bitrev}(i)} \mod q$) sont déployés comme bytecode d'un contrat ``donnée''. \texttt{extcodecopy} est l'instruction EVM pour lire le bytecode d'un contrat externe. C'est \textbf{plus gas-efficace} que stocker les données en \texttt{storage} (SLOAD = 2100 gas) ou les passer en calldata. Le bytecode est accessible en lecture seule via \texttt{EXTCODECOPY} avec seulement quelques gas.

\textbf{2. L'assembleur Yul dans le butterfly :}\\
Le butterfly de Cooley-Tukey est :
\begin{lstlisting}[style=solidity]
assembly {
    // t = mulmod(zeta, b, q)
    let t := mulmod(mload(ptr_b), zeta, Q)
    // a, b = a + t, a - t (mod q)
    mstore(ptr_b, addmod(sub(Q, t), mload(ptr_a), Q))
    mstore(ptr_a, addmod(t, mload(ptr_a), Q))
}
\end{lstlisting}
\texttt{mulmod/addmod} sont des opcodes EVM natifs qui effectuent l'opération modulo en une seule instruction sans overflow.

\textbf{3. Nombre de twiddle factors :}\\
128 valeurs (= n/2 = 256/2), codées sur 32 bytes chacune (uint256), donc $128 \times 32 = 4096$ bytes de bytecode.

\textbf{4. Coût gas NTT :}\\
Environ 200K--400K gas pour une NTT complète. C'est pour cela que ETHFALCON (Falcon) est visé : sa NTT sur des anneaux plus petits coûte bien moins.
\end{tcolorbox}
\end{solution}

%──────────────────────────────────────────────────────────────────────
\subsection{Exercice P10 : Écrire un test Foundry}
%──────────────────────────────────────────────────────────────────────

\begin{exercice}{P10 --- Ton premier test Forge}
Crée un test Foundry minimal qui :
\begin{enumerate}
  \item Génère une paire de clés en Python
  \item Crée une signature en Python
  \item Appelle le contrat Solidity \texttt{verify()} et vérifie le résultat
\end{enumerate}

\textbf{Fichier à créer :} \texttt{ETHDILITHIUM-main/test/MyTest.t.sol}
\end{exercice}

\begin{solution}
\begin{lstlisting}[style=python]
# Étape 1 : Générer les données de test en Python
# python-ref/gen_test_data.py
from Falcon_dilithium.Falcon_dilithium import Dilithium2
import json

msg = b"SuVa Protocol v1"
pk, sk = Dilithium2.keygen()
sig    = Dilithium2.sign(sk, msg)
valid  = Dilithium2.verify(pk, msg, sig)

assert valid, "Signature invalide!"

data = {
    "pk":  pk.hex(),
    "sk":  sk.hex(),
    "sig": sig.hex(),
    "msg": msg.hex(),
}
with open("test_data.json", "w") as f:
    json.dump(data, f, indent=2)
print("Test data saved to test_data.json")
print(f"pk={len(pk)}B, sk={len(sk)}B, sig={len(sig)}B")
\end{lstlisting}

\begin{lstlisting}[style=solidity]
// ETHDILITHIUM-main/test/MyTest.t.sol
// SPDX-License-Identifier: MIT
pragma solidity ^0.8.20;

import "forge-std/Test.sol";
// Importer le contrat de vérification (ajuster le chemin)
import "../src/ZKNOX_dilithium2.sol";

contract MyDilithiumTest is Test {
    ZKNOX_dilithium2 verifier;

    function setUp() public {
        verifier = new ZKNOX_dilithium2();
    }

    function test_verify_simple() public view {
        // Ces bytes doivent être générés par le script Python ci-dessus
        // et copiés ici (ou lus depuis un fichier avec vm.readFile)
        bytes memory pk  = hex"<COLLER_PK_HEX_ICI>";
        bytes memory sig = hex"<COLLER_SIG_HEX_ICI>";
        bytes memory msg = hex"53755661_50726f74_6f636f6c_763100"; // "SuVa Protocol v1"

        bool result = verifier.verify(pk, msg, sig);
        assertTrue(result, "La verification Dilithium doit reussir");
    }

    function test_reject_bad_message() public view {
        bytes memory pk      = hex"<COLLER_PK_HEX_ICI>";
        bytes memory sig     = hex"<COLLER_SIG_HEX_ICI>";
        bytes memory bad_msg = hex"00000000"; // Message incorrect

        bool result = verifier.verify(pk, bad_msg, sig);
        assertFalse(result, "La verification doit echouer avec un mauvais message");
    }
}
\end{lstlisting}

\begin{lstlisting}[style=bash]
# Lancer les tests
cd ETHDILITHIUM-main
forge test --match-contract MyDilithiumTest -vvv --gas-report
\end{lstlisting}
\end{solution}

\newpage
%══════════════════════════════════════════════════════════════════════
\section{Récapitulatif : Checklist de maîtrise}
%══════════════════════════════════════════════════════════════════════

\begin{tcolorbox}[colback=success_bg, colframe=sol_green, title=\textbf{Checklist --- Tu es prêt quand tu peux :}]
\begin{itemize}[label=$\square$]
  \item[\checkmark] Lancer \texttt{pytest} et voir 4 tests passer
  \item[\checkmark] Écrire un script Python keygen/sign/verify sans copier-coller
  \item[\checkmark] Expliquer pourquoi le rejection sampling boucle
  \item[\checkmark] Dire la différence entre Dilithium2 et ETHDilithium2
  \item[\checkmark] Implémenter \texttt{decompose()} de mémoire
  \item[\checkmark] Savoir ce que stocke \texttt{extcodecopy} dans ZKNOX\_NTT.sol
  \item[\checkmark] Écrire un test Foundry minimal qui appelle \texttt{verify()}
  \item[\checkmark] Modifier un paramètre (eta, gamma1) et expliquer l'impact
\end{itemize}
\end{tcolorbox}

\begin{tcolorbox}[colback=warn_bg, colframe=orange!80!black, title=\textbf{Prochaines étapes}]
\begin{enumerate}
  \item \textbf{Intégrer Falcon (ETHFALCON)} : Le ZKNox a déjà une implémentation. Il faut l'adapter au projet SuVa.
  \item \textbf{Développer le Vault Contract} : Le contrat Solidity qui wrap USDT/USDC en sUSDT/sUSDC.
  \item \textbf{ERC-4337 Wallet} : Le smart contract wallet qui utilise Dilithium/Falcon comme validateur de signature.
  \item \textbf{SPHINCS+} : Fallback hash-based (pas encore commencé, ZKNox n'a pas d'implémentation EVM).
\end{enumerate}
\end{tcolorbox}

\vspace{1cm}
\begin{center}
  {\color{suva_blue}\rule{0.5\linewidth}{1pt}}\\[0.3cm]
  {\small\color{gray} Document créé par \textbf{Mohamed Lamine MBENGUE} ---
  \href{https://portfolio-alamine.web.app}{\texttt{portfolio-alamine.web.app}}}
\end{center}

\end{document}
